\section{The Life for Which We Were Made}

Heaven is eternal life with God: the perfect communion of the blessed with the Father, through the Son, in the Holy Spirit. Those who die in God's grace and friendship, once wholly purified, see him face to face and live with Christ forever (CCC 1023--1029). God himself is their happiness. The good for which the Christian hopes is therefore greater than relief from suffering, reunion with the beloved dead, or the continuation of a life already known. All these desires find their proper place within the promise that we shall belong wholly to God and possess him without end.

Before his Passion, Jesus speaks of the Father's house and promises to receive his disciples to himself (Jn 14:1--6). The destination has the character of a home because the Father receives his children there through the Son. Christ is the way. He goes before us through his death and Resurrection and desires that where he is, we also may be. Preparation for heaven begins by receiving that desire of Christ as the truth about our future, even when our own desire remains hesitant or divided.

The same farewell discourse joins the future dwelling to God's present indwelling. The Father and the Son make their home with the one who loves Christ and keeps his word; the Spirit is promised as the abiding Paraclete (Jn 14:15--26). Heaven brings to fullness the communion that grace begins. God already draws the baptized person into friendship, teaches him to pray, forgives him, and gives him the charity by which he can love. Yet this beginning is lived in faith, amid weakness and struggle. The fullness remains a gift to be received.

In the prayer of John 17, Jesus names eternal life as knowing the Father and the Son whom he has sent. This knowledge involves the whole person in the love in which the Son glorifies the Father. Jesus asks that his disciples be with him to behold his glory, and he speaks of the Father's love for him before the foundation of the world (Jn 17:1--5, 24--26). Our salvation opens into a love that did not begin with us and will never be exhausted by us.

The Father, Son, and Holy Spirit are the one God whose own life is freely shared with the redeemed. We receive this communion by grace and remain the creatures God has made and loves. The Son who became man brings us to his Father; the Spirit makes us living members of Christ. Every creaturely joy in heaven is received within that communion, whose source and measure remain God himself.

In the light of Christ's promise, we can examine whether the projects death will interrupt are serving the life he offers. Some responsibilities become more urgent: fidelity to those entrusted to us, repentance for our sins, the work of reconciliation. Other ambitions lose their claim to rule the heart. We begin to distinguish what can be offered to God from what we have been trying to possess apart from him.

Such preparation has room for affection, gratitude, and sorrow. The thought of heaven does not make the death of someone loved painless. Nor must we pretend that leaving our present work and companionship costs nothing. Christ promises a good that exceeds our present knowledge, and we approach it with a heart that still knows separation from within. Hope enables that heart to entrust itself to him. We follow the Savior who knows us, who died for us, and who wills to gather his people into his glory.

The Four Last Things lead toward this fulfillment. Death makes our pilgrimage finite; judgment discloses its truth; hell reveals the terrible finality of refusing God. Heaven names the end for which mercy calls us, the good toward which conversion moves. A soul can become so accustomed to considering only its faults that it forgets what forgiveness is restoring. God heals us for friendship with himself. Remembering the last things each year gives us occasion to consider what we have allowed to weaken that friendship and to seek its renewal.

